\documentclass[11pt,letterpaper,oneside]{report}
\usepackage[T1]{fontenc}
\usepackage[utf8]{inputenc}
\usepackage{lmodern}
\usepackage[margin=0.90in,headheight=15pt]{geometry}
\usepackage{microtype}
\usepackage[svgnames,table]{xcolor}
\usepackage{amsmath,amssymb,booktabs,longtable,tabularx,array}
\usepackage{enumitem}
\usepackage[most]{tcolorbox}
\usepackage{listings}
\usepackage{fancyhdr}
\usepackage{titlesec}
\usepackage{hyperref}
\usepackage{xurl}
\usepackage{lastpage}
\definecolor{ink}{HTML}{14273C}
\definecolor{primary}{HTML}{0B5379}
\definecolor{accent}{HTML}{B66A30}
\definecolor{pale}{HTML}{F1F5F8}
\definecolor{codebg}{HTML}{F5F7F8}
\hypersetup{colorlinks=true,linkcolor=primary,urlcolor=primary,pdftitle={Linux and UNIX System Programming Engineering Suite}}
\titleformat{\chapter}[display]{\normalfont\LARGE\bfseries\color{ink}}{\color{accent}\large\MakeUppercase{\chaptername}\ \thechapter}{10pt}{}[\vspace{0.4ex}\titlerule]
\titleformat{\section}{\Large\bfseries\color{primary}}{\thesection}{0.7em}{}
\titleformat{\subsection}{\large\bfseries\color{ink}}{\thesubsection}{0.7em}{}
\setlist[itemize]{leftmargin=*,topsep=3pt,itemsep=2pt}
\setlist[enumerate]{leftmargin=*,topsep=3pt,itemsep=3pt}
\setlength{\parskip}{5pt}
\setlength{\parindent}{0pt}
\lstdefinestyle{sysc}{language=C,basicstyle=\ttfamily\footnotesize,backgroundcolor=\color{codebg},keywordstyle=\bfseries\color{primary},commentstyle=\color{ForestGreen},stringstyle=\color{Brown},breaklines=true,showstringspaces=false,tabsize=4,frame=single,rulecolor=\color{LightSteelBlue},columns=fullflexible,keepspaces=true}
\lstset{style=sysc}
\newtcolorbox{engineeringnote}[1][]{colback=pale,colframe=primary,boxrule=0.55pt,arc=1.5mm,title=Engineering note,fonttitle=\bfseries,#1}
\newtcolorbox{modernnote}[1][]{colback=white,colframe=accent,boxrule=0.55pt,arc=1.5mm,title=Modern extension --- not in the 2010 textbook,fonttitle=\bfseries,#1}
\newcommand{\sys}[1]{\texttt{#1}}
\pagestyle{fancy}\fancyhf{}\fancyhead[L]{\small\color{ink} Linux / UNIX Engineering Suite}\fancyhead[R]{\small\color{primary}\nouppercase{\leftmark}}\fancyfoot[C]{\small\thepage\ / \pageref{LastPage}}
\renewcommand{\headrulewidth}{0.45pt}
\setcounter{tocdepth}{1}
\setlength{\emergencystretch}{2em}
\begin{document}
\begin{titlepage}\centering\vspace*{2cm}{\color{primary}\Large ENGINEERING REFERENCE SERIES\par}\vspace{1.25cm}{\color{ink}\Huge\bfseries Prerequisite-based Curriculum\par}\vspace{0.6cm}{\Large 24-week engineer training path, milestones and evaluations\par}\vspace{2cm}\rule{0.75\textwidth}{1.2pt}\par\vspace{0.55cm}{\large An original companion to Michael Kerrisk\textquotesingle s 2010 text}\par\vspace{0.3cm}{\large Edition 1.0\quad |\quad October 2026}\par\vfill{\small Document 09\quad |\quad pdfLaTeX-compatible standalone source}\par\end{titlepage}
\tableofcontents\clearpage
\chapter{Audience, outcomes and operating model}
This 24-week curriculum is designed for C-capable software engineers who can compile a program and use a shell but may be new to Linux/UNIX system interfaces. Learners progress from descriptor and memory correctness to concurrency, IPC, network servers, security and modern Linux APIs. The attached textbook supplies the core topic sequence; contemporary tools and standards are identified as additions.

\section{Exit competencies}
\begin{itemize}
\item Design and defend a POSIX resource ownership and cleanup contract.
\item Build reproducible C interfaces with checked errors and bounded processing.
\item Reason about fork/exec, signals, pthread synchronization and memory mapping.
\item Select IPC and socket designs using threat, portability and workload constraints.
\item Implement an unprivileged daemon with controlled shutdown, observability and rigorous failure testing.
\item Identify differences among POSIX.1-2008-era material, POSIX.1-2024 and Linux-specific facilities.
\end{itemize}
\chapter{Prerequisite dependency graph}
\begin{center}\fcolorbox{primary}{pale}{\parbox{.88\linewidth}{\centering
C basics + shell + compiler warnings\\$\downarrow$\\
I/O contracts + descriptors + memory + runtime identity\\$\downarrow$\\
Process lifecycle + signals + threads + synchronization\\$\downarrow$\\
IPC + socket protocols + bounded server architectures\\$\downarrow$\\
Security boundaries + modern Linux optional APIs\\$\downarrow$\\
Capstone: resilient least-privilege service}}
\end{center}\section{Advance criteria}
\begin{enumerate}
\item After week 3: demonstrate binary-safe I/O with cleanup and correct allocation boundaries.
\item After week 9: demonstrate race-aware child reaping and signal-safe termination.
\item After week 14: demonstrate deadlock-conscious pthread logic, least-privilege design and safe plugin loading.
\item After week 21: implement bounded local and remote protocols with predictable EOF and backpressure.
\item After week 24: deliver capstone passing the security and reliability rubric.
\end{enumerate}
\chapter{Weekly schedule and reading traceability}
\footnotesize\setlength{\tabcolsep}{3pt}\begin{longtable}{p{.065\linewidth}p{.245\linewidth}p{.115\linewidth}p{.12\linewidth}p{.35\linewidth}}\toprule
Week & Technical focus & Book chapters & Labs & Weekly evidence\\\midrule\endfirsthead
\toprule Wk & Technical focus & Chapters & Labs & Evidence\\\midrule\endhead
1 & C pointers, builds, shell, test discipline & 1--3 & 01, 02 & Compiler and API provenance inventory\\
2 & Descriptor semantics and binary I/O & 4--5 & 02, 03 & Checked binary copy with failure tests\\
3 & Process layouts and allocation & 6--7 & 04 & Memory-safe vector and address-space diagram\\
4 & Identity, credentials and runtime limits & 8--12 & 05 & Identity inspector and bounds worksheet\\
5 & Buffering and filesystem invariants & 13--15 & 02, 05 & Durability and permissions review\\
6 & Extended metadata and path traversal & 16--19 & 06 & Directory-fd and symlink-race report\\
7 & Signal semantics & 20--22 & 08 & Signal safety and managed shutdown\\
8 & Clocks, timeouts, spawning & 23--25 & 07, 09 & Deadline loop and controlled child launch\\
9 & Child reaping and exec transitions & 26--28 & 09 & Multi-child supervisor with status capture\\
10 & Pthreads and guarded state & 29--30 & 10 & Bounded queue correctness proof\\
11 & Thread safety and cancellation & 31--33 & 10 & Cancellation cleanup and stress traces\\
12 & Job control and resource limits & 34--36 & 09 & Resource-limited process test\\
13 & Services, secrets and privileges & 37--39 & 05, 20 & Threat model and least-privilege controls\\
14 & Runtime linking and plugins & 40--42 & 11 & Versioned plugin with loader allowlist\\
15 & Pipe and FIFO lifecycles & 43--44 & 12 & Streaming IPC prototype\\
16 & System V IPC & 45--48 & 13 & Object lifetime and queue exhaustion study\\
17 & Memory maps and POSIX IPC & 49--54 & 14 & Versioned shared-memory protocol\\
18 & File locking and local sockets & 55--57 & 15 & Local peer authorization design\\
19 & TCP/IP and name resolution & 58--59 & 16 & Dual-stack connection client\\
20 & Server architectures & 60--61 & 17 & Bounded TCP server with framing\\
21 & Terminals and readiness I/O & 62--64 & 17, 18 & epoll draining and PTY recovery\\
22 & Modern Linux optional facilities & Modern & 19 & API probes and compatibility decision record\\
23 & Security and observability integration & Review all & 20 & Capstone integration tests and traces\\
24 & Capstone acceptance and presentation & Review all & 20 & Release package, threat model, performance report\\
\bottomrule\end{longtable}\normalsize\setlength{\tabcolsep}{6pt}\chapter{Competency tracks and alternative routes}
\section{Systems foundations}
Prerequisites: Weeks 1--9. Emphasis: File and memory APIs; reliable process execution; build and debug tools.

\section{Concurrent service engineering}
Prerequisites: Systems foundations plus Weeks 10--18. Emphasis: Pthreads, signal strategies, queues, IPC protocols, scheduling and recovery.

\section{Network service engineering}
Prerequisites: Systems foundations plus Weeks 15--21. Emphasis: UNIX-domain and TCP/IP socket protocols, framing, readiness and lifecycle controls.

\section{Security specialization}
Prerequisites: Systems foundations plus Weeks 13--24. Emphasis: Pathname races, credential transitions, restricted execution, privilege separation, diagnostics and system hardening.

\chapter{Capstone specification}
\section{Project: local-control, loopback service}
Implement a small, unprivileged service that accepts a bounded versioned request over a loopback TCP socket. Provide a separate UNIX-domain administration channel that can query health and request a graceful shutdown. Use a finite work queue, explicit thread ownership, controlled child processes only if needed, and a local state file created in a trusted directory.

\section{Required functional contracts}
\begin{enumerate}
\item Requests: version and payload length validated before memory allocation; maximum frame 64 KiB.
\item Concurrency: fixed worker count, bounded queue, and backpressure behavior defined.
\item Network: dual-stack design review; loopback-only binding for the training build.
\item Shutdown: SIGTERM creates a controlled drain; no async-unsafe handler activity.
\item Files: fd-scoped directory operations with an explicit symlink and failure policy.
\item Privileges: no setuid binary or ambient root requirement; record process credentials at startup.
\item Observability: request counters, queue depth, failure codes and p95 response time without payload leaks.
\item Testing: malformed lengths, connection churn, premature EOF, bind conflict, unavailable resources and SIGTERM during work.
\end{enumerate}
\section{Deliverable package}
\begin{itemize}
\item Architecture decision record and trust-boundary diagram
\item C source, Makefile and explicit compiler flags
\item Unit, integration and negative test scripts
\item POSIX/Linux extension classification sheet
\item Risk register and threat model
\item Performance run including throughput, p50/p95/p99 and memory peak
\item Operator runbook covering startup, shutdown, logs, and known limitations
\end{itemize}
\chapter{Assessment, mentoring and spaced review}
\begin{longtable}{p{.23\linewidth}p{.60\linewidth}}\toprule
Review cadence & Review activity\\\midrule\endhead
Every lab & Compiler warning-free build and paired error-path review\\
Every 2 weeks & Live whiteboard: ownership, state machine and one failure scenario\\
Weeks 3, 9, 14, 21 & Checkpoint with resubmission allowed after corrective work\\
Week 24 & Capstone demonstration, documentation defense and triage simulation\\
\bottomrule\end{longtable}Assessment weighting mirrors the laboratory rubric: 30\% behavior/tests, 20\% error handling, 20\% security/trust, 15\% lifecycle/concurrency, 10\% portability and 5\% clarity. A passing program must correctly terminate, release owned objects and reject malformed inputs; test volume alone cannot offset critical bugs.

\chapter{Recommended reference update plan}
\begin{itemize}
\item At each new course cohort, compare covered interfaces with current POSIX.1-2024 texts and Linux manual pages.
\item Re-check minimum kernel and libc version for Linux-only lab features; avoid fragile distribution-specific assumptions.
\item Rotate compiler and sanitizer versions and audit false positives/false negatives.
\item Retain reference outputs from at least one supported low-resource virtual machine and one contemporary Linux host.
\end{itemize}

\chapter*{Source and currency notes}
\addcontentsline{toc}{chapter}{Source and currency notes}
\textbf{Foundational source:} Michael Kerrisk, \emph{The Linux Programming Interface: A Linux and UNIX System Programming Handbook}, No Starch Press, 2010, ISBN 978-1-59327-220-3. The textbook's chapter titles and start pages organize the study guide; explanations, tasks, and code in this suite are independently authored educational material, not reproductions of the original text or its example programs.

\textbf{Current reference baseline:} The Open Group / IEEE, \emph{POSIX.1-2024, Issue 8}, \url{https://pubs.opengroup.org/onlinepubs/9799919799/}; Linux man-pages project, \url{https://man7.org/linux/man-pages/}; Linux kernel documentation, \url{https://docs.kernel.org/}. For Linux-only extensions consult the relevant per-interface manual pages and kernel documentation for the target kernel, architecture and libc; availability can differ materially.

\textbf{Scope:} This is an engineering study and implementation companion. It does not claim to restate a normative standard or guarantee a specific Linux kernel's behavior. System calls, libc wrappers and higher-level library functions are explicitly distinguished where material.
\end{document}
